\subsubsection{Detectie op basis van machine learning}
\label{subsubsec:machinelearning-detectie}

\emph{Machine learning (ML)} vervangt de handmatig ingestelde gewichten door modellen \autocite{Salloum2022}. Deze modellen worden getraind op échte phishingmails en leren zelf de correlatie tussen kenmerken die frequent voorkomen bij phishingmails. De regels zelf verdwijnen dus niet, enkel de manier waarop ze worden gecombineerd verandert.

\paragraph{Kenmerkselectie.}
Bij kenmerkselectie wordt getracht enkel de kenmerken te behouden die bijdragen aan het onderscheiden tussen phishingmails en legitieme e-mails. Op die manier wordt het model sneller en minder gevoelig voor ruis. \textcite{Guyon2003} onderscheiden drie families van methoden om relevante kenmerken te selecteren:

\begin{enumerate}
    \item \textbf{Filtermethoden.} Bij filtermethoden wordt elk kenmerk afzonderlijk beoordeeld, nog voordat het model wordt getraind. Een veelgebruikte maatstaf is \emph{information gain}, die meet in welke mate de waarde van een kenmerk de onzekerheid over de klasse (phishing of legitiem) vermindert \autocite{Forman2003}. Kenmerken die laag scoren, worden verwijderd. Filtermethoden zijn snel en onafhankelijk van het gekozen model, maar houden geen rekening met de samenhang tussen kenmerken. Twee kenmerken die afzonderlijk zwak scoren, kunnen samen toch een belangrijke meerwaarde bieden \autocite{Guyon2003}.
    \item \textbf{Wrappermethoden.} Bij wrappermethoden wordt het model getraind op verschillende subsets van kenmerken, waarna de subset met de beste prestatie wordt behouden \autocite{Kohavi1997}. Op die manier wordt er wel rekening gehouden met de samenhang tussen kenmerken. Het nadeel is de sterk oplopende rekenkost, want het aantal mogelijke subsets groeit exponentieel met het aantal kenmerken, en het model moet tegen elke subset getraind worden \autocite{Kohavi1997}.
    \item \textbf{Embedded methoden.} Bij embedded methoden gebeurt de selectie tijdens het trainen van het model zelf \autocite{Guyon2003}. Een voorbeeld is L1-regula\-risatie. Daarbij kent het model aan elk kenmerk een gewicht toe dat aangeeft in welke mate dat kenmerk meetelt in de uiteindelijke beslissing. De som van de absolute waarden van alle gewichten wordt beperkt, zodat het model maar over een beperkt budget beschikt en niet elk kenmerk zwaar kan laten doorwegen. Het model moet dus kiezen aan welke kenmerken het gewicht wordt besteed. Kenmerken die weinig bijdragen aan de classificatie krijgen een gewicht van exact nul en spelen dan geen rol meer in het model \autocite{Fatih2024}. Hoe streng dit budget is, wordt bepaald door een instelbare parameter, waarbij geldt dat een strenger budget leidt tot minder kenmerken.
\end{enumerate}

\paragraph{Veelgebruikte modellen.}
Voor phishingdetectie kunnen verschillende modellen gebruikt worden:

\begin{enumerate}
    \item \textbf{Gradient boosting.} Gradient boosting is geschikt voor gestructureerde, tabulaire data, zoals de header-, authenticatie- en URL-ken\-merken uit Sectie~\ref{subsec:kenmerken-detectie}. Het bouwt een reeks kleine beslissingsbomen, waarbij elke nieuwe boom probeert de fouten van de voorgaande bomen te corrigeren. Een beslissingsboom bestaat uit opeenvolgende vragen over de kenmerken van een e-mail. Op basis van de antwoorden wordt het bericht naar een classificatie geleid. Onderzoek van \textcite{Grinsztajn2022} toont aan dat boomgebaseerde modellen, zoals gradient boosting, op tabulaire data vaak competitief blijven met \emph{deep learning}-methoden, terwijl ze minder rekenkracht vereisen. XGBoost is een veelgebruikte implementatie van gradient boosting \autocite{Chen2016}.
    \item \textbf{Transformermodels.} Transformermodels zijn het meest geschikt voor ongestructureerde data, zoals de tekstinhoud van een bericht. Modellen zoals BERT worden eerst op grote hoeveelheden tekst voorgetraind en daarna bijgetraind op als phishing bestempelde e-mails \autocite{Devlin2019}. Ze houden rekening met de volledige context van een zin. 
    \item \textbf{Large Language Models (LLM).} \emph{LLM's} kunnen ook zonder taakspecifieke training worden ingezet. \textcite{Koide2025} ontwikkelden een systeem waarbij inhoud van een e-mail als prompt aan een LLM werd aangeboden, waarna het model een beoordeling teruggaf.
\end{enumerate}

De hierboven besproken modellen leveren een antwoord op de vraag of een e-mail phishing is of niet. Echter blijken veel classifiers volgens \textcite{Niculescu-Mizil2005} geen kansen te produceren die goed gekalibreerd zijn, waardoor een extra kalibratiestap nodig is.

\paragraph{Logistische regressie.}
\label{subsubsec:logistische-regressie}
Een model dat voor phishingdetectie goede prestaties behaalt en bovendien interpreteerbaar is, is logistische regressie \autocite{Alhuzali2025}. Elk kenmerk krijgt daarbij een gewicht dat aangeeft hoeveel het bijdraagt aan de score. Dat is vergelijkbaar met de gewogen scores bij rulebasede systemen (zie Sectie~\ref{subsubsec:rulebased-detectie}), maar met gewichten die uit data worden geleerd.

\paragraph{Nadelen.}
\label{subsubsec:machinelearning-detectie:nadelen}
Machine learning lost het probleem van handmatig ingestelde gewichten op, maar brengt eigen nadelen met zich mee. Een eerste nadeel is de afhankelijkheid van data. Een model leert enkel wat in de trainingsdata aanwezig is, waardoor de kwaliteit rechtstreeks afhangt van de kwaliteit en de representativiteit ervan \autocite{Mendes2025}.

Een tweede nadeel is dat phishingaanvallen voortdurend evolueren. Hierdoor vertonen nieuwe berichten andere kenmerken dan de berichten uit de trainingsdata. \textcite{Fernando2026} beschrijven dit als \emph{concept drift}. Het gevolg hiervan is dat het model geleidelijk minder goed gaat presteren, omdat het op oude data is getraind.

Tot slot komen er voor LLM's nog twee specifieke nadelen bovenom. De verklaring die een LLM bij zijn beoordeling geeft, hoeft niet de werkelijke reden van de beslissing weer te geven, waardoor ze overtuigend kan klinken zonder correct te zijn \autocite{Turpin2023}. Verder kan een aanvaller via \emph{indirect prompt injection} instructies in de e-mailinhoud verwerken om de beoordeling van het model te beïnvloeden \autocite{Greshake2023}.